\section{逆问题框架（Inverse Problems）}

第二类引导技术的核心思想是将条件生成建模为\textbf{逆问题}（Inverse Problem）。这是一个更加结构化和通用的框架，适用于更广泛的应用场景。

\subsection{逆问题的定义}

\begin{importantbox}{逆问题的形式化定义}
给定：
\begin{itemize}
    \item 一个\textbf{前向映射函数} $\mathcal{A}$，将干净数据 $x$ 映射到部分观测 $y$：$y = \mathcal{A}(x) + n$，其中 $n$ 是可能的观测噪声。
    \item 一个\textbf{观测值} $y$。
\end{itemize}
目标：找到一个合理的 $x$，使得 $\mathcal{A}(x) \approx y$。

由于 $\mathcal{A}$ 通常\textbf{不可逆}（non-invertible），给定 $y$ 可能存在多个合理的 $x$，因此这本质上是一个\textbf{条件生成}问题。
\end{importantbox}

\subsection{逆问题的典型应用}

逆问题在图像处理中有大量应用，也可扩展到其他数据模态：

\begin{table}[H]
\centering
\begin{tabular}{lll}
\toprule
\textbf{任务} & \textbf{前向映射 $\mathcal{A}$} & \textbf{说明} \\
\midrule
高斯去模糊 & 高斯卷积核 & 从模糊图像恢复清晰图像 \\
超分辨率 & 下采样操作 & 从低分辨率恢复高分辨率 \\
图像修复 & 掩码操作 & 从部分观测恢复完整图像 \\
CT/MRI 重建 & 投影操作 & 从多角度 2D 投影重建 3D \\
图像着色 & 去色操作 & 从灰度图恢复彩色图 \\
\bottomrule
\end{tabular}
\caption{逆问题的典型应用场景}
\end{table}

\begin{knowledgebox}{逆问题的广泛性}
逆问题不仅存在于图像领域。在分子设计、蛋白质折叠、机器人规划等领域，任何``从部分观测推断完整信息''的问题都可以建模为逆问题。如果你能为某个新领域定义合理的前向映射 $\mathcal{A}$，就可以利用本讲介绍的技术来解决。
\end{knowledgebox}

\subsection{本章小结}

逆问题提供了一个统一的框架，将推理时引导生成与经典的信号处理/计算成像问题联系起来。关键要素是：已知的前向映射 $\mathcal{A}$、观测值 $y$、以及预训练的无条件生成模型。接下来我们推导如何利用贝叶斯公式将逆问题求解转化为对分数函数的修改。


